\documentclass[12pt]{article}
\usepackage[utf8]{inputenc}
\usepackage[T1]{fontenc}
\usepackage[spanish]{babel}
\usepackage{graphicx}
\usepackage{tabularx}
\usepackage{hyperref}
\usepackage{geometry}
\usepackage[table]{xcolor}
\usepackage{fancyhdr}
\usepackage{titlesec}

% Márgenes amplios y limpios
\geometry{
  top=2.5cm,
  bottom=2.5cm,
  left=3cm,
  right=3cm
}

% Colores institucionales
\definecolor{ipnGuinda}{RGB}{108, 29, 69}
\definecolor{tablaHeader}{RGB}{230, 230, 230}

% Formato de títulos
\titleformat{\section}
  {\normalfont\Large\bfseries\color{ipnGuinda}}{\thesection}{1em}{}
\titleformat{\subsection}
  {\normalfont\large\bfseries\color{ipnGuinda}}{\thesubsection}{1em}{}

% Configuración de encabezados y pies de página
\pagestyle{fancy}
\fancyhf{}
\fancyhead[L]{\textcolor{ipnGuinda}{\textbf{Análisis de Negocio}}}
\fancyhead[R]{\small Módulo de Protección NNA}
\fancyfoot[C]{\thepage}
\renewcommand{\headrulewidth}{0.5pt}
\renewcommand{\headrule}{\hbox to\headwidth{\color{ipnGuinda}\leaders\hrule height \headrulewidth\hfill}}

% Definición de comandos tipo CDT mejorados
\newcommand{\cdtInstrucciones}[1]{}
\newcommand{\cdtEmpty}{\vspace{-1em}}

\newenvironment{Usuario}[2]{
    \vspace{1.5em}\noindent\textcolor{ipnGuinda}{\textbf{\Large #1}}\par\noindent\textit{#2}\par\vspace{0.8em}
    \begin{description}
}{
    \end{description}
}

\newenvironment{cdtEntidad}[2]{
    \vspace{1.5em}\noindent\textcolor{ipnGuinda}{\textbf{\Large Entidad: #1 (\hypertarget{#1}{#2})}}\par\vspace{0.8em}
    \renewcommand{\arraystretch}{1.3}
    \tabularx{\textwidth}{|l|l|X|c|}
    \hline
    \rowcolor{tablaHeader}
    \textbf{Atributo} & \textbf{Tipo} & \textbf{Descripción} & \textbf{Req.} \\
    \hline
}{
    \endtabularx\vspace{1.5em}
}

\newcommand{\brAttr}[5]{
    #2 & #3 & #4 & #5 \\
    \hline
}

\newcommand{\cdtEntityRelSection}{
    \rowcolor{tablaHeader}
    \multicolumn{4}{|c|}{\textbf{Relaciones}} \\
    \hline
}

\newcommand{\brRel}[3]{
    \multicolumn{2}{|l|}{\textbf{#1}} & \multicolumn{2}{l|}{#3} \\
    \hline
}

\newcommand{\brRelComposition}{Composición}
\newcommand{\brRelAgregation}{Agregación}
\newcommand{\brRelGeneralization}{Generalización}
\newcommand{\brRelParticipation}{Asociación}

\newcommand{\CasoDeUso}[5]{
    \vspace{1.5em}\noindent\textcolor{ipnGuinda}{\textbf{\Large CU-#1: #2}}\par\vspace{0.8em}
    \renewcommand{\arraystretch}{1.3}
    \begin{tabular}{|>{\columncolor{tablaHeader}\bfseries}p{3.5cm}|p{11.5cm}|}
    \hline
    Actor Principal & #3 \\
    \hline
    Descripción & #4 \\
    \hline
    Flujo Principal & #5 \\
    \hline
    \end{tabular}\vspace{1.5em}
}

\begin{document}
\thispagestyle{empty} % sin número de página en la portada

% --------- ENCABEZADO CON LOGOS ----------
\begin{minipage}{0.49\linewidth}
  \includegraphics[width=3cm]{C:/Users/damia/OneDrive - Instituto Politecnico Nacional/Imagenes/IPN.png}
\end{minipage}
\hfill
\begin{minipage}{0.49\linewidth}
  \flushright
  \includegraphics[width=3cm]{C:/Users/damia/OneDrive - Instituto Politecnico Nacional/Imagenes/ESCOM.png}
\end{minipage}


% --------- CUERPO CENTRADO VERTICALMENTE ----------
\begin{center}
\vspace*{1cm} % espacio después de logos

{\Large \textbf{Instituto Politécnico Nacional}} \\[0.2cm]
{\large \textbf{Escuela Superior de Cómputo}} \\[1cm]

{\large \textbf{Materia:}} \\[0.2cm]
{\Large ANÁLISIS Y DISEÑO DE SISTEMAS} \\[1cm]

{\Huge \textbf{Análisis de Negocio\\ Módulo de Protección NNA y Familiograma}} \\[0.7cm]
\textcolor{ipnGuinda}{\rule{0.6\linewidth}{1.5pt}} \\[1cm]

{\large \textbf{Profesor:}} \\[0.2cm]
{\Large  Ulises Vélez Saldaña} \\[1cm]

{\large \textbf{Alumno:}} \\[0.2cm]
\begin{flushleft}
\hspace{4cm} % sangría estética
\begin{tabular}{ll}
Canales Zendreros Diego Damian  \\ 
Rueda Manzano, Jennifer \\
\end{tabular}
\end{flushleft}

{\large \textbf{Grupo:}} \\[0.2cm]
{\Large 4BV1} \\[1cm]

{\large Ciudad de México, 12 de Mayo de 2026}

\vfill % empuja hacia abajo el último bloque si sobra espacio
\end{center}

\newpage
\tableofcontents
\newpage

\section{Introducción}
Este documento describe el modelo de negocio para el Módulo de Protección de Niñas, Niños y Adolescentes (NNA) y la herramienta de Familiograma dentro del sistema Xolix 3.0. Su propósito es definir los actores involucrados, los términos clave del negocio, el modelo de dominio del problema y las principales reglas de negocio aplicables a esta iteración.

\section{Actores del sistema}

\begin{Usuario}{\hypertarget{ADirector}{\subsection{Director}}}{
        Máxima autoridad de la fundación, responsable de la supervisión global.
    }
    \item[Responsabilidades:] \cdtEmpty
    \begin{itemize}
        \item Supervisar todos los casos NNA activos e inactivos.
        \item Eliminar casos NNA o registros del sistema en situaciones excepcionales.
        \item Monitorear los planes de acción globales.
    \end{itemize}
    \item[Perfil:] \cdtEmpty
    \begin{itemize}
        \item Profesional en gestión de instituciones sociales o psicología clínica.
    \end{itemize}
    \item[Procesos en los que participa:] Gestión general de expedientes, Supervisión de Casos NNA.
\end{Usuario}

\begin{Usuario}{\hypertarget{ACoordinador}{\subsection{Coordinador / Psicólogo}}}{
        Personal operativo especializado responsable del trato directo con los NNA y sus familias.
    }
    \item[Responsabilidades:] \cdtEmpty
    \begin{itemize}
        \item Registrar nuevos casos NNA.
        \item Realizar y registrar entrevistas familiares.
        \item Construir y mantener actualizado el familiograma del caso.
        \item Registrar observaciones no verbales durante las sesiones.
        \item Generar planes de acción basados en las entrevistas.
    \end{itemize}
    \item[Perfil:] \cdtEmpty
    \begin{itemize}
        \item Licenciatura en Psicología, Trabajo Social o afines.
        \item Experiencia en trato con menores vulnerados.
    \end{itemize}
    \item[Procesos en los que participa:] Entrevistas de diagnóstico, Creación de Familiogramas, Planes de Acción.
\end{Usuario}

\section{Términos del Negocio}

\begin{description}
    \item[\hypertarget{tNNA}{NNA:}] Acrónimo de Niña, Niño o Adolescente. Sujeto principal de protección y atención de la fundación.
    \item[\hypertarget{tCaso}{Caso NNA:}] Expediente integral que agrupa toda la información de un NNA específico, incluyendo sus datos personales, familiares, psicológicos y legales.
    \item[\hypertarget{tFamiliograma}{Familiograma:}] Representación gráfica (grafo) interactiva del núcleo y la red familiar de un NNA. Muestra de manera visual a los individuos y la naturaleza de las relaciones entre ellos.
    \item[\hypertarget{tPersonaFamiliar}{Persona Familiar:}] Individuo que pertenece al entorno (familiar, social o de cuidado) del NNA, registrado en el sistema con sus datos sociodemográficos y rol.
    \item[\hypertarget{tRelacionFamiliar}{Relación Familiar:}] Vínculo específico entre dos Personas Familiares dentro de un caso. Puede ser de diversos tipos (biológica, legal, conflictiva, protectora, etc.) y tener direccionalidad.
    \item[\hypertarget{tPlanAccion}{Plan de Acción:}] Conjunto de tareas y procesos generados (a menudo automáticamente según el nivel de negación) para dar seguimiento a un Caso NNA.
\end{description}

\section{Modelo del dominio del problema}

A continuación se describen las entidades principales que conforman el modelo de datos del Módulo de Protección NNA y Familiograma.

\begin{cdtEntidad}{CasoNNA}{Caso NNA}
    \brAttr{nna\_nombre}{Nombre del NNA}{Cadena}{Nombre completo del NNA}{Sí}
    \brAttr{nna\_edad}{Edad}{Entero}{Edad en años del NNA}{No}
    \brAttr{nna\_genero}{Género}{Enum}{Masculino, Femenino, No Binario, Otro}{No}
    \brAttr{estado}{Estado}{Enum}{Activo o Cerrado}{Sí}
    \cdtEntityRelSection
    \brRel{\brRelComposition}{PersonaFamiliar}{Un \hyperlink{CasoNNA}{Caso} contiene muchas \hyperlink{PersonaFamiliar}{Personas}}
    \brRel{\brRelComposition}{RelacionFamiliar}{Un \hyperlink{CasoNNA}{Caso} contiene muchas \hyperlink{RelacionFamiliar}{Relaciones}}
    \brRel{\brRelComposition}{Familiograma}{Un \hyperlink{CasoNNA}{Caso} tiene un \hyperlink{Familiograma}{Familiograma}}
\end{cdtEntidad}

\begin{cdtEntidad}{PersonaFamiliar}{Persona Familiar}
    \brAttr{nombre}{Nombre}{Cadena}{Nombre de la persona}{Sí}
    \brAttr{rol\_en\_familia}{Rol}{Cadena}{Ej. Padre, Madre, Abuelo}{No}
    \brAttr{tipo\_simbolo}{Símbolo}{Enum}{Normal, NNA Clave, Agresor, Cuidador, Fallecido}{Sí}
    \brAttr{vive\_con\_nna}{Vive con NNA}{Booleano}{Indica si cohabita con el NNA}{Sí}
    \brAttr{es\_responsable\_legal}{Responsable}{Booleano}{Indica si tiene la tutela legal}{Sí}
    \cdtEntityRelSection
    \brRel{\brRelAgregation}{CasoNNA}{Pertenece a un \hyperlink{CasoNNA}{Caso NNA}}
\end{cdtEntidad}

\begin{cdtEntidad}{RelacionFamiliar}{Relación Familiar}
    \brAttr{tipo\_relacion}{Tipo}{Enum}{Biológica, Legal, Conflictiva, Protectora, etc.}{Sí}
    \brAttr{descripcion}{Descripción}{Cadena}{Notas adicionales sobre el vínculo}{No}
    \brAttr{bidireccional}{Bidireccional}{Booleano}{Si el vínculo es mutuo}{Sí}
    \cdtEntityRelSection
    \brRel{\brRelParticipation}{PersonaFamiliar}{Une una \hyperlink{PersonaFamiliar}{Persona Origen} con una Persona Destino}
    \brRel{\brRelAgregation}{CasoNNA}{Pertenece a un \hyperlink{CasoNNA}{Caso NNA}}
\end{cdtEntidad}

\begin{cdtEntidad}{Familiograma}{Familiograma}
    \brAttr{grafo\_json}{Grafo JSON}{JSON}{Estructura visual de nodos y aristas para el editor ReactFlow}{Sí}
    \brAttr{imagen\_url}{Imagen URL}{Cadena}{Snapshot de la última versión gráfica}{No}
    \cdtEntityRelSection
    \brRel{\brRelComposition}{HistorialFamiliograma}{Genera múltiples \hyperlink{HistorialFamiliograma}{Versiones en el Historial}}
\end{cdtEntidad}

\begin{cdtEntidad}{HistorialFamiliograma}{Historial Familiograma}
    \brAttr{version}{Versión}{Entero}{Número de versión autoincremental}{Sí}
    \brAttr{grafo\_json}{Grafo}{JSON}{Copia del estado del grafo en esta versión}{Sí}
    \brAttr{notas\_version}{Notas}{Cadena}{Comentario de los cambios (ej. "Restaurado desde v2")}{No}
    \cdtEntityRelSection
    \brRel{\brRelAgregation}{Familiograma}{Pertenece a un \hyperlink{Familiograma}{Familiograma}}
\end{cdtEntidad}

\section{Reglas de Negocio}

\begin{description}
    \item[\textbf{RN-NNA-01 Unicidad de relaciones:}] Una relación familiar origen-destino debe asociar a dos personas diferentes (una persona no puede tener una relación consigo misma).
    \item[\textbf{RN-NNA-02 Versionado automático:}] Cada vez que se actualiza el \hyperlink{tFamiliograma}{Familiograma} (el JSON del grafo cambia), el sistema debe guardar automáticamente una copia del estado anterior en el \hyperlink{tHistorialFamiliograma}{Historial} para permitir restauraciones.
    \item[\textbf{RN-NNA-03 Dependencia de personas para relaciones:}] Para poder crear una \hyperlink{tRelacionFamiliar}{Relación Familiar} en el sistema, deben existir al menos dos \hyperlink{tPersonaFamiliar}{Personas} registradas en el mismo Caso.
    \item[\textbf{RN-NNA-04 Acceso a casos:}] Un coordinador solo puede modificar los casos NNA de los que es creador o responsable asignado. El director tiene acceso global a todos los casos.
\end{description}

\section{Casos de Uso}

A continuación, se describen los 10 casos de uso principales correspondientes a la gestión de NNA y el ciclo de vida del familiograma interactivo.

\CasoDeUso{01}{Registrar Caso NNA}
{Coordinador / Psicólogo}
{Permite al actor registrar a un nuevo menor de edad en situación de vulnerabilidad abriendo un expediente.}
{
\begin{enumerate}
    \item El actor selecciona "Nuevo Caso NNA" en el dashboard.
    \item El sistema muestra el formulario de datos generales.
    \item El actor ingresa nombre, CURP, edad, género y motivo de ingreso.
    \item El sistema valida la información y crea el registro del caso en estado "Activo".
\end{enumerate}
}

\CasoDeUso{02}{Consultar Casos NNA}
{Director, Coordinador / Psicólogo}
{Permite visualizar la lista de expedientes activos e inactivos, con herramientas de filtrado.}
{
\begin{enumerate}
    \item El actor accede al Módulo de Protección NNA.
    \item El sistema lista los casos NNA recientes.
    \item El actor utiliza los filtros (estado, rango de edad, género) o la barra de búsqueda.
    \item El sistema actualiza la lista mostrando coincidencias.
\end{enumerate}
}

\CasoDeUso{03}{Registrar Persona Familiar}
{Coordinador / Psicólogo}
{Añade a un individuo (padre, madre, hermano, cuidador, agresor) al expediente de un NNA para después vincularlo al familiograma.}
{
\begin{enumerate}
    \item El actor ingresa a un Caso NNA y va a la pestaña "Personas Familiares".
    \item El actor hace clic en "Agregar Persona".
    \item El actor completa los datos sociodemográficos y define el rol o símbolo clínico.
    \item El sistema guarda a la persona y la hace disponible para el familiograma.
\end{enumerate}
}

\CasoDeUso{04}{Registrar Relación Familiar}
{Coordinador / Psicólogo}
{Crea un vínculo clínico o descriptivo entre dos personas familiares registradas en un mismo caso.}
{
\begin{enumerate}
    \item El actor accede a la pestaña "Relaciones Familiares" del caso.
    \item El actor selecciona a la "Persona Origen" y a la "Persona Destino".
    \item El actor escoge el tipo de relación (biológica, conflictiva, distante, etc.).
    \item El sistema registra la relación bidireccional o unidireccional según aplique.
\end{enumerate}
}

\CasoDeUso{05}{Modificar Grafo del Familiograma}
{Coordinador / Psicólogo}
{Permite utilizar el editor visual (React Flow) para acomodar los nodos y trazar las líneas de las relaciones familiares de forma gráfica.}
{
\begin{enumerate}
    \item El actor ingresa al "Editor de Familiograma".
    \item El sistema carga los nodos (personas) y aristas (relaciones) guardadas previamente.
    \item El actor arrastra los nodos para acomodarlos o crea nuevas aristas visualmente.
    \item El actor presiona "Guardar".
    \item El sistema actualiza el grafo JSON y genera un \textbf{Snapshot en el Historial}.
\end{enumerate}
}

\CasoDeUso{06}{Consultar Historial del Familiograma}
{Coordinador / Psicólogo}
{Visualizar los estados pasados (versiones) de un familiograma a lo largo del tiempo.}
{
\begin{enumerate}
    \item El actor accede a "Historial del Familiograma".
    \item El sistema despliega una línea de tiempo con las versiones guardadas.
    \item El actor selecciona una versión antigua.
    \item El sistema dibuja el grafo exacto correspondiente a esa versión temporal.
\end{enumerate}
}

\CasoDeUso{07}{Restaurar Versión del Familiograma}
{Coordinador / Psicólogo}
{Revierte el estado actual del familiograma a un snapshot histórico (ej. por cambios erróneos).}
{
\begin{enumerate}
    \item El actor visualiza una versión antigua en el "Historial del Familiograma".
    \item El actor hace clic en "Restaurar esta Versión".
    \item El sistema clona el JSON de esa versión antigua y lo establece como la versión actual (nueva versión $N+1$).
    \item El sistema recarga el editor visual con la estructura restaurada.
\end{enumerate}
}

\CasoDeUso{08}{Exportar Reporte del Caso NNA}
{Director, Coordinador / Psicólogo}
{Genera un documento consolidado (JSON/PDF) con los datos del menor, sus familiares, relaciones y el diagrama visual.}
{
\begin{enumerate}
    \item El actor ingresa a "Resumen Integral" del caso.
    \item El actor hace clic en el botón "Exportar / Imprimir".
    \item El sistema consolida la información de las 4 tablas (Caso, Personas, Relaciones, Familiograma).
    \item El sistema descarga el archivo en el dispositivo del usuario.
\end{enumerate}
}

\CasoDeUso{09}{Registrar Entrevista Familiar (Wizard)}
{Coordinador / Psicólogo}
{Asiste al psicólogo para registrar las observaciones de una entrevista en una secuencia guiada paso a paso.}
{
\begin{enumerate}
    \item El actor ingresa a "Nueva Entrevista" desde el hub del caso.
    \item El sistema muestra un Wizard interactivo (Paso 1: Motivo, Paso 2: Observaciones, Paso 3: Familiograma dinámico).
    \item El actor completa cada sección y avanza.
    \item Al finalizar, el sistema guarda la sesión de entrevista y actualiza el expediente.
\end{enumerate}
}

\CasoDeUso{10}{Cerrar Caso NNA}
{Director}
{Cambia el estado de un expediente a inactivo cuando el menor egresa o concluye su proceso legal.}
{
\begin{enumerate}
    \item El actor (con rol Director) ingresa a un Caso NNA.
    \item En la configuración del caso, selecciona "Cerrar Expediente".
    \item El sistema solicita una confirmación y un motivo de cierre.
    \item El actor confirma.
    \item El sistema bloquea futuras modificaciones al familiograma y cambia el estado a "Cerrado".
\end{enumerate}
}

\end{document}
